% !TeX root = ../nc-quadric-avatar.tex
% 2026-11-23 / v1.0
\section{相対avatarの構造頂点と二つの行列式開}
\label{sec:qa-relative-vertex-determinants}

本節は，体上の全構造頂点判定を，任意の非被約底変更と両立する
有限表示の開条件へ移す．まず相対曲線と実際の線束だけから
二つの乗法行列を作る．この幾何適合開の構成と，実際の相対
pointed $\HH$族の認識は区別し，後者に必要な入力を末尾に明記する．

\subsection{相対乗法束と全基底変更}

$T$を$k$-schemeとし，$f:\mathcal D\to T$をproper，flat，
finite presentationな射とする．$\mathcal F,\mathcal G,\mathcal P$を
$\mathcal D$上の可逆層とし，各幾何点$\bar t\to T$において
\[
 (\mathcal D_{\bar t},\mathcal F_{\bar t},
   \mathcal G_{\bar t},\mathcal P_{\bar t})
\]
が本稿のsquare曲線型avatarであると仮定する．
曲線の被約性も$T$の被約性も仮定しない．
とくに$\mathcal P_{\bar t}\in\Pic^0(\mathcal D_{\bar t})
\setminus\{\OO\}$である．

\begin{lemma}[二つのrank四乗法束]
\label{lem:qa-relative-pencil-multiplication-bundles}
$\mathcal V(L)=f_*L$と書く．
\[
\mathcal V(\mathcal F),\quad\mathcal V(\mathcal G),\quad
\mathcal V(\mathcal F\mathcal P^{-1}),\quad
\mathcal V(\mathcal G\mathcal P^{-1})
\]
は階数二，$\mathcal W=\mathcal V(\mathcal F\mathcal G\mathcal P^{-1})$
は階数四の有限局所自由層である．
これらの高次直像は零で，形成は任意の$T'\to T$と可換する．
従って切断の乗法はrank四束間の二つの写像
\begin{align}
 \mu_G:&\ \mathcal V(\mathcal F)\otimes
          \mathcal V(\mathcal G\mathcal P^{-1})\longrightarrow\mathcal W,
          \label{eq:qa-relative-mu-g}\\
 \mu_F:&\ \mathcal V(\mathcal G)\otimes
          \mathcal V(\mathcal F\mathcal P^{-1})\longrightarrow\mathcal W
          \label{eq:qa-relative-mu-f}
\end{align}
を定め，両写像も任意の底変更と可換する．
全幾何fiberにおける階数は三または四である．
\end{lemma}
\begin{proof}
補題~\ref{lem:qa-core-curve-vanishing}より，各fiberで対象線束の
$H^1$は零，$H^0$の次元は記載した二または四である．
各$Rf_*L$はperfectで，任意の導来底変更と可換する
（\href{https://stacks.math.columbia.edu/tag/0B91}
{Stacks Project, Tag 0B91, Lemma 36.30.4}）．
その全幾何fiberが次数零に集中するので，有限自由複体の可逆minorを
順に消去すれば，$Rf_*L$自体が記載した階数の局所自由層となる．
従って導来底変更は通常の底変更となる．
これは非被約底上でコホモロジー次元が一定という理由だけから
局所自由性を推論するものではない．
乗法の自然性は導来底変更およびprojection formulaと両立し，
階数三または四という主張は
補題~\ref{lem:qa-as10-rank-wall-pencils}から従う．
\end{proof}

\subsection{幾何適合開とschemeとしての悪いlocus}

\begin{theorem}[標準構造頂点の有限表示開部分関手]
\label{thm:qa-relative-vertex-determinant-open}
体上のAS/avatar対応で各幾何fiberから得られる$i$番標準構造頂点について，
$T$上のschemeの部分関手$\mathscr U_i$を，すべての幾何fiberの
指定pointingが強い$\HH$に属する底変更$T'\to T$の全体とする．
これは次の有限表示open immersion $U_i\hookrightarrow T$で表される：
\begin{equation}
 U_0=U_3=T,\qquad
 U_1=D(\det\mu_F),\qquad U_2=D(\det\mu_G).
 \label{eq:qa-relative-square-vertex-open}
\end{equation}
行列式はそれぞれ
$\det\mathcal W\otimes\det(\operatorname{dom}\mu)^{-1}$の切断である．
すべての$T'\to T$について
\[
 U_i(\mathcal D_{T'},\mathcal F_{T'},\mathcal G_{T'},\mathcal P_{T'})
 =U_i\times_T T'
\]
がschemeとして成立する．

悪い閉部分schemeを
\[
 Z_F=V\bigl(\operatorname{Fitt}_0\coker\mu_G\bigr),\qquad
 Z_G=V\bigl(\operatorname{Fitt}_0\coker\mu_F\bigr)
\]
と置けば，$U_1=T\setminus Z_G$，$U_2=T\setminus Z_F$である．
$Z_F,Z_G$も有限表示であり，任意の底変更と可換する．
各写像は局所的に
\begin{equation}
 \operatorname{diag}(1,1,1,\delta)
 \label{eq:qa-relative-wall-schur-normal-form}
\end{equation}
と同値であり，対応する悪い閉部分は$(\delta)$で定義される．
その余核は悪い閉部分上の線束のpushforwardである．
\end{theorem}
\begin{proof}
前補題により乗法束と行列式はすべての底変更と可換する．
定理~\ref{thm:qa-general-square-pointings}によって，幾何fiberの
一番頂点は$\mu_F$が可逆の場合に，二番頂点は$\mu_G$が可逆の場合に
ちょうど強い$\HH$へ属する．
線束の切断が全幾何点で非零であることは，その切断が可逆であることと
同値であり，非被約底でも成り立つ．
従って各部分関手は記載した行列式の可逆開で表される．
線束を局所自明化すればprincipal openであるため，
open immersionは有限表示である．

rank四束間の写像の第零Fitting idealは行列式で生成される．
Fitting idealの形成は任意の環準同型と可換する
（\href{https://stacks.math.columbia.edu/tag/07Z6}
{Stacks Project, Tag 07Z6, Lemma 15.8.4(3)--(4)}）．
全幾何fiberで階数が三以上なので，三次minorの可逆開が$T$を覆う．
その上で行列の消去を行えば
\eqref{eq:qa-relative-wall-schur-normal-form}を得る．
余核は局所的に$\OO_T/(\delta)$であり，これらの局所表示は
$V(\delta)$上の可逆層として貼り合う．
ここで$\delta$が非零因子であるとは仮定していない．
\end{proof}

\begin{corollary}[Rotated族と最小十本矢族の開部分]
\label{cor:qa-relative-rotated-as10-open}
rotated族から
$\mathcal F_s=\mathcal F_r$，
$\mathcal G_s=\mathcal H_r\mathcal F_r^{-1}\mathcal P^{-1}$
によって二乗法を作ると，rotatedの各位相の幾何適合開は
\[
 U^{\rm rsq}_0=D(\det\mu_F),\quad
 U^{\rm rsq}_1=D(\det\mu_G),\quad
 U^{\rm rsq}_2=U^{\rm rsq}_3=T.
\]
全幾何fiberが最小十本矢型である族では，偶数位相の幾何適合開は空，
奇数位相の幾何適合開は$T$である．
これらはいずれも任意の底変更と可換する．
\end{corollary}
\begin{proof}
前者は$R_i=E_{i+1}$と相対line bundleの式に，後者は
系~\ref{cor:qa-general-rotated-as10-pointings}に従う．
rotatedから得た線束の各fiberはsquare avatarなので，
前補題と前定理をそのまま適用できる．
\end{proof}

\subsection{普遍楕円パラメータ上のCartier壁}

\begin{proposition}[普遍壁はCartier，任意引戻しではFitting閉部分]
\label{prop:qa-relative-universal-cartier-walls}
$E\to S$を標数零の滑らかな楕円曲線族，$F$を相対次数二線束とし，
$J=\Pic^0(E/S)$と置く．
\[
 V=(J\setminus0)\times_S(J\setminus0),\qquad
 (P,\eta)\in V,\quad G=F\eta
\]
上で，必要ならPoincare束を持つfppf chartへ移って乗法束を作る．
このときscheme-theoreticに
\[
 Z_F=\{\eta=P\},\qquad Z_G=\{\eta=P^{-1}\}
\]
であり，両者は$V$上の有効Cartier因子である．
両者の交わりは$J[2]\setminus0$で，$S$上有限etale次数三である．
相対余次元二で横断的に交わる．

任意のavatar族への引戻しについては，悪い閉部分は上記Fitting理想の
引戻しであるが，有効Cartier因子であるとは限らない．
引き戻した局所行列式が非零因子である場合に限り，Cartierである．
\end{proposition}
\begin{proof}
次数二pencil $A$と$B$について，相対評価列は
\[
 0\to f^*\det\mathcal V(A)\otimes BA^{-1}
 \to f^*\mathcal V(A)\otimes B\to AB\to0
\]
を与える．正次数項の高次直像消滅から
\[
 Rf_*(BA^{-1})\otimes\det\mathcal V(A)
 \simeq
 [\mathcal V(A)\otimes\mathcal V(B)\xrightarrow{\mu}\mathcal V(AB)]
\]
を得る．右辺の二項は次数零，一に置く．
従って行列式の消滅は，次数零線束$BA^{-1}$のcohomology jumpを
そのscheme構造も含めて記録する．

普遍$J$の零断面上では対象線束は基底からの線束を除き自明であり，
kernelとcokernelは各一の線束となる．
ここでは$Rf_*\OO_E$のunitからの分裂と
$f_*\OO_E=\OO_S$を用いる
（\href{https://stacks.math.columbia.edu/tag/0E6B}
{Stacks Project, Tag 0E6B, Lemma 75.26.7}）．
したがって行列式は零断面上でscheme-theoreticに零となる．
各幾何fiberで零断面に垂直な非零接方向
$\xi\in H^1(E_s,\OO_{E_s})$を取ると，
対応する一次変形$Q_\epsilon$の接続写像は
\[
 H^0(\OO_{E_s})\longrightarrow H^1(\OO_{E_s}),\qquad1\longmapsto\xi
\]
である．従って残る一行一列の微分は非零で，行列式は零断面で一次の零点を持つ．
零断面の局所Cartier方程式を$z$とすれば，行列式は$zu$と書け，
このfiber上の計算から$u$は単元である．
ゆえにこの同定は非被約な$S$上でもscheme-theoreticに成立する．

$\mu_G$では$BA^{-1}=\eta P^{-1}$，
$\mu_F$では$BA^{-1}=\eta^{-1}P^{-1}$であるから，
記載した二つのgraphを得る．
滑らかな相対一次元群$J$におけるdiagonalおよび反転のgraphは
有効Cartierである．交点では$P^2=\OO$，かつ$P\ne\OO$であり，
$2$は可逆なので$J[2]\setminus0$は有限etale次数三である．
二つの相対法方向は$d\eta-dP$，$d\eta+dP$であり，
行列式$2$が可逆であることから横断性が従う．
Poincare束の基底からのline bundle twistは行列式の局所方程式を
単元倍するだけなので，これらの閉部分はchartから降下する．
最後の主張は局所正規形の$\delta$がCartier方程式となる条件そのものである．
\end{proof}

\begin{example}[非被約底での階数三条件]
普遍壁に横断する一次変形を
$T=\Spec k[\epsilon]/(\epsilon^2)$へ引き戻すと，局所乗法行列は
\[
 \operatorname{diag}(1,1,1,\epsilon)
\]
となる．全幾何fiberの階数は三だが，行列式は$T$上恒等的に零ではない．
悪いFitting閉部分は$V(\epsilon)$で，有効Cartier因子ではない．
一方，対応する幾何適合開は空である．
したがって幾何fiberの階数だけからscheme-theoreticな壁への因数分解，
kernel束の局所自由性，またはCartier性を推論しない．
\end{example}

\subsection{相対pointed認識へ進むための有限入力}

\begin{lemma}[既存H2開への相対因数分解]
\label{lem:qa-relative-foundation-open-factorization}
$T$を本稿の相対モジュライsite内の基底とする．
実際の有限局所自由foundation代数族と全積から
射$\lambda:T\to\mathfrak F_0$が与えられ，
その全幾何fiberが実現可能開$\mathfrak F_{\HH}$に入るとする．
このとき$\lambda$は一意に$\mathfrak F_{\HH}$を経由し，
定理~\ref{thm:qa-intrinsic-h2-open-in-foundation-stack}により
実際の相対pointed $\HH$族を定める．
この構成は任意のsite内の底変更と同型射に適合する．
\end{lemma}
\begin{proof}
開部分への因数分解は位相空間の像で判定できるので，
非被約底についても$\lambda$は$\mathfrak F_{\HH}$を経由する．
同定$\mathfrak F_{\HH}\simeq\mathfrak H_2^{\rm rul}$を合成すればよい．
開埋込とこの同値の自然性から底変更・同型射への適合性が従う．
\end{proof}

この補題で必要なのは，四頂点の全Hom束と全合成を持つ
\emph{実際の}相対foundation族，$\mathfrak F_0$の条件，
および幾何fiberにおけるnormal foundationとしての認識である．
行列式開の上にあるという事実だけからこれらを省略しない．
さらに，こうして得た$\HH$族と，別に与えた相対AS代数の$qgr$ heartを
同一視するには，構造対象を保つ相対tilting比較，Serre完成と
相対AZ再構成を確認する必要がある．
本節の二行列式だけで全AS/QP stackを$\HHstack$と同定したとは主張しない．
